\chapter{Introduction}

Electrochemistry is a branch of chemistry that is concerned with the
relationship between electrical applied potential and chemical
change. These reactions arise at the interface between an electronic
conductor and an ionic conductor, where charge transfer drives
oxidation and reduction of chemical species. The study of these
reactions has applications in biochemical sensing, energy storage,
technology and many others that underpin much of modern biology and
nanotechnology. In particular,
electroanalytical techniques have enabled the detection of trace
analytes in biological fluids, the development of rechargeable
batteries and fuel cells, and the fabrication of nanoscale electronic
devices, underlining the broad relevance of electrochemical science
across disciplines.

Voltammetry is a common technique to analyse an electroactive
species in solution. Voltammetry measures the current that is induced
into an electrochemical circuit as the potential of the working
electrode is varied with respect to the reference electrode. The
current response arises from Faradaic processes, in which electrons
are transferred between the electrode surface and the electroactive
species, and its magnitude reflects both the kinetics of this charge
transfer and the rate at which species are transported to the
electrode. Voltammetric techniques are distinguished by how the
applied potential is varied over time. In cyclic voltammetry, the
potential is swept linearly between two limits and then reversed,
producing a characteristic current--potential curve from which
thermodynamic and kinetic parameters may be inferred. Other
waveforms, such as the sinusoidal perturbations of electrochemical
impedance spectroscopy or the staircase profiles of square-wave
voltammetry, offer complementary sensitivity and resolution.

Voltammetry has been extensively studied mathematically and
numerically. The electron transfer kinetics at the electrode surface
are described by the Butler--Volmer model, which relates the rate of
the electrode reaction to the applied overpotential through
exponential anodic and cathodic contributions, or alternatively by
the Marcus--Hush framework, which grounds the kinetics in a
quantum-mechanical treatment of solvent reorganisation. At
equilibrium, the Nernst equation dictates the ratio of oxidised to
reduced species concentrations at the electrode surface as a function
of the applied potential. The transport of electroactive species
through the solution is governed by Fick's laws of diffusion, which,
combined with the electrode boundary conditions described above,
yield a partial differential equation (PDE) that fully characterises
the evolution of the concentration profile in the cell. The solution
of this system--the so-called forward problem--predicts the
current response given a known set of physical parameters, and has
been extensively studied both analytically and numerically
\cite{compton2014understanding}.

We are interested in the inverse of this process: given the
measured current response from a voltammetric experiment, we wish to
infer the underlying physical parameters that characterise the
electroactive species in solution. These parameters may include the
formal potential, diffusion coefficients, and electron transfer rate
constants, knowledge of which provides quantitative insight into the
thermodynamics and kinetics of the redox system under study. This
constitutes the inverse problem, and its solution is considerably
more challenging than the forward problem. The mapping from
parameters to current is nonlinear and the measurements are
corrupted by noise, meaning that naive inversion is typically ill-posed
and that uncertainty in the inferred parameters must be carefully
quantified.

Many different techniques have been employed to solve this inverse
problem. Deterministic optimisation methods, such as least-squares
fitting, provide point estimates of the parameters but offer no
principled account of uncertainty. Bayesian inference addresses this
limitation by treating the parameters as random variables and seeking
their posterior distribution conditioned on the observed data. Markov
chain Monte Carlo (MCMC) methods have emerged as the dominant
computational approach to posterior sampling in electrochemistry,
enabling rigorous uncertainty quantification without requiring
analytical tractability of the posterior \cite{gavaghan2018bayesian}. However,
standard MCMC algorithms such as the random-walk Metropolis--Hastings
sampler can mix poorly in high-dimensional or geometrically complex
parameter spaces, and the requirement to repeatedly solve the forward
PDE makes the burn-in phase computationally expensive.

With the advent of modern automatic differentiation (AD) software,
it is now possible to efficiently compute exact derivatives of
arbitrary numerical programs, including complete PDE solvers, with
respect to their inputs \cite{chen2026differentiable}. This capability opens
new avenues for gradient-based inference in electrochemistry. In
particular, the gradient of the simulated current with respect to the
model parameters can be obtained at a cost comparable to a single
forward solve, enabling the use of gradient-informed algorithms both
during the burn-in phase and throughout the sampling procedure. During
burn-in, gradient information can be exploited to drive the chain
rapidly toward regions of high posterior probability, dramatically
reducing the number of forward solves required before convergence.
Once in the posterior, gradient-based samplers such as Hamiltonian Monte
Carlo (HMC) use the gradient to make more informed proposals,
resulting in higher acceptance rates and faster exploration of the
posterior landscape compared with gradient-free alternatives.

This dissertation demonstrates the applicability of gradient-based
methodologies to solve the inverse problem in electrochemical
experiments. We develop a differentiable PDE solver for common
voltammetric waveforms, validate it against known analytical results,
and integrate it with gradient-informed MCMC algorithms to perform
Bayesian inference of electrochemical parameters. We show that the
use of automatic differentiation yields substantial improvements in
computational efficiency over gradient-free methods, and we
illustrate the approach on both synthetic and experimental data.
